\section{2011 Geometry \& Topology}

\subsection{Individual Exam}

\begin{exercise}
\label{geometry:2011:1}
Suppose $M$ is a closed smooth $n$-manifold.

(a) Does there always exist a smooth map $f : M \to S^{n}$ from $M$ into the $n$-sphere, such that $f$ is essential (i.e., $f$ is not homotopic to a constant map)? Justify your answer.

(b) Same question, replacing $S^{n}$ by the $n$-torus $T^{n}$.
\end{exercise}

\begin{solution}
\textbf{Geometric Intuition:} The $n$-sphere $S^n$ can be thought of as the result of collapsing the boundary of an $n$-disk to a point. For any $n$-manifold, we can always find a small local $n$-disk and "wrap" it around the sphere, which should capture some non-trivial topological information. In contrast, the $n$-torus $T^n$ has a very different "flat" topology, and its universal cover is contractible, which severely limits the maps from simpler manifolds like $S^n$ to it.

(a) \textbf{Yes.} For any closed connected $n$-manifold $M$, we can construct an essential map $f: M \to S^n$ as follows:
Pick a point $p \in M$ and a small neighborhood $U$ homeomorphic to an open $n$-disk $D^n$. Let $K = M \setminus U$. Since $M$ is a manifold, $K$ is a closed subset. We can define a map $f: M \to M/K$\footnote{Intuitively, $M \setminus K$ is an open $n$-disk. Collapsing the exterior $K$ to a point is topologically equivalent to taking a closed disk $D^n$ and identifying its entire boundary $\partial D^n$ to a single point, which precisely forms the $n$-sphere $S^n$.}. The quotient space $M/K$ is homeomorphic to $D^n / \partial D^n$, which is the $n$-sphere $S^n$.
If $M$ is oriented, this map $f$ has degree 1, which means it induces an isomorphism $f_*: H_n(M; \mathbb{Z}) \to H_n(S^n; \mathbb{Z}) \cong \mathbb{Z}$. Since $f_* \neq 0$, $f$ cannot be null-homotopic.
If $M$ is non-orientable, we consider homology with $\mathbb{Z}_2$ coefficients. The map $f$ induces an isomorphism $f_*: H_n(M; \mathbb{Z}_2) \to H_n(S^n; \mathbb{Z}_2) \cong \mathbb{Z}_2$. Again, $f_* \neq 0$ implies $f$ is essential.

(b) \textbf{No.} Consider $M = S^n$ for $n > 1$. Any map $f: S^n \to T^n$ is homotopic to a constant map.
The $n$-torus $T^n$ is the product of $n$ circles $(S^1)^n$. Its universal cover is $\mathbb{R}^n$, which is contractible. A basic result in algebraic topology states that any map from a simply connected space (like $S^n$ for $n > 1$) to a space with a contractible universal cover is null-homotopic. Formally, $\pi_n(T^n) \cong \pi_n(S^1) \times \dots \times \pi_n(S^1) = 0$ for $n > 1$. Thus, every map $f: S^n \to T^n$ represents the zero element in $\pi_n(T^n)$\sidenote{A map $f: S^n \to X$ is said to \textbf{represent} an element of the homotopy group $\pi_n(X)$ because $\pi_n(X)$ consists of \textit{classes} of maps under the relation of homotopy. If $f$ represents the 'zero element' (the identity), it means it belongs to the same class as the constant map, implying it can be continuously shrunk to a point.}, making it null-homotopic.
\end{solution}

\begin{exercise}
\label{geometry:2011:2}
Suppose $(X, d)$ is a compact metric space and $f : X \to X$ is a map so that $d(f(x), f(y)) = d(x, y)$ for all $x, y$ in $X$. Show that $f$ is an onto map.
\end{exercise}

\begin{solution}
\textbf{Geometric Intuition:} An isometry preserves all distances. In a "finite" world (like a finite set), an injective map is always surjective. Compactness is the topological generalization of finiteness. If the map were not onto, we could keep "moving away" from the missed region, but in a compact space, we eventually run out of room and must come back near where we started, contradicting the distance-preserving property.

\textbf{Proof:} Suppose $f$ is not onto. Then there exists some $y \in X$ such that $y \notin f(X)$. Since $f$ is an isometry, it is continuous. The image $f(X)$ of a compact set $X$ under a continuous map is compact, and hence closed. Since $y \notin f(X)$, there exists some $\epsilon > 0$ such that the ball $B(y, \epsilon)$ is disjoint from $f(X)$.
Consider the sequence defined by $y_0 = y$ and $y_{n+1} = f(y_n)$ for $n \geq 0$. Note that $y_n = f^n(y)$. Since $f$ is an isometry, for any $n > m \geq 0$:
\[ d(y_n, y_m) = d(f^m(y_{n-m}), f^m(y_0)) = d(y_{n-m}, y_0) \]
By construction, $y_{n-m} = f(y_{n-m-1}) \in f(X)$. Since $y_0 = y \notin f(X)$ and the distance from $y$ to $f(X)$ is at least $\epsilon$, we have $d(y_{n-m}, y_0) \geq \epsilon$.
Thus, $d(y_n, y_m) \geq \epsilon$ for all $n \neq m$. This means the sequence $\{y_n\}$ has no convergent subsequence, as no two terms are closer than $\epsilon$. However, this contradicts the assumption that $X$ is compact (in a compact metric space, every sequence must have a convergent subsequence). Therefore, $f$ must be onto.
\end{solution}

\begin{exercise}
\label{geometry:2011:3}
Let $C_{1}, C_{2}$ be two linked circles\sidenote{The \textbf{linking number} $Lk(C_1, C_2)$ of two disjoint oriented curves $C_1, C_2 \subset \mathbb{R}^3$ is a topological invariant that measures how they are interlinked. It can be defined as the degree of the map $f: C_1 \times C_2 \to S^2$ given by $f(u, v) = \frac{u-v}{|u-v|}$. Analytically, it is given by the \textbf{Gauss linking integral}: $Lk(C_1, C_2) = \frac{1}{4\pi} \oint_{C_1} \oint_{C_2} \frac{\mathbf{r}_1 - \mathbf{r}_2}{|\mathbf{r}_1 - \mathbf{r}_2|^3} \cdot (d\mathbf{r}_1 \times d\mathbf{r}_2)$. A standard example is the \textbf{Hopf link}, consisting of two circles where each passes through the disk bounded by the other (e.g., $x^2+y^2=1, z=0$ and $(x-1)^2+z^2=1, y=0$), resulting in $Lk = \pm 1$.} in $\mathbb{R}^{3}$. Show that $C_{1}$ cannot be homotopic to a point in $\mathbb{R}^{3} \setminus C_{2}$.
\end{exercise}



\begin{solution}
\textbf{Prerequisites:}
\begin{itemize}
    \item Homotopy and fundamental groups.
    \item Concept of Linking Number.
    \item Degree of a continuous map.\sidenote{The \textbf{degree} of a map $f: M \to N$ between oriented $n$-manifolds measures how many times $M$ 'wraps around' $N$. Formally, if $[M]$ and $[N]$ are fundamental classes, it is the integer $d$ such that $f_*([M]) = d[N]$. For smooth maps, it is the sum of the signs of the Jacobian at the preimages of a regular value.}
\end{itemize}

\textbf{Intuition:}
Imagine two interlocking physical rings. Being "linked" is a topological property that prevents one ring from being continuously shrunk to a point or pulled far away without passing exactly through the other ring. 

If you could continuously deform the first ring $C_1$ to a single point without ever touching the second ring $C_2$, it would imply that $C_1$ was never truly "wrapped around" $C_2$ to begin with. We make this mathematically rigorous using a tool called the linking number, which remains completely unchanged under continuous deformations. If $C_1$ shrinks to a point, the linking number would be zero, contradicting the fact that the rings are linked.

\textbf{Steps:}
\begin{enumerate}
    \item We formalize the "linked" property using the \textbf{Linking Number}, denoted $Lk(C_1, C_2)$. For two disjoint closed curves in $\mathbb{R}^3$, the linking number is a topological invariant. It can be defined as the degree of the map $\Phi: C_1 \times C_2 \to S^2$ given by:
    \[ \Phi(u, v) = \frac{u - v}{|u - v|} \! \]
    Since the circles are linked, we know $Lk(C_1, C_2) \neq 0$.

    \item Now, suppose for the sake of contradiction that $C_1$ is homotopic to a point in the space $\mathbb{R}^3 \setminus C_2$. This means there is a continuous deformation (homotopy) $H: S^1 \times [0,1] \to \mathbb{R}^3 \setminus C_2$.\sidenote{A circle $C \subset \mathbb{R}^3$ is formally defined as the image of a continuous map (an embedding) from the unit circle $S^1$. Therefore, properties like 'homotopy to a point' are analyzed using the parameterizing map $\gamma: S^1 \to \mathbb{R}^3$.}

    \item In this homotopy, $H(s, 0)$ is the standard parameterization of $C_1$, and $H(s, 1)$ is a single constant point $p$ that does not lie on $C_2$. Importantly, the homotopy never intersects $C_2$ at any time.

    \item This homotopy $H$ induces a continuous deformation of our map $\Phi$. Specifically, $\Phi$ is homotopic to a new map $\Psi(s, t) = \frac{p - v(t)}{|p - v(t)|}$, where $v(t)$ parameterizes $C_2$. 

    \item Notice that the map $\Psi$ only depends on the parameter $t$ of the second curve, and is completely independent of the parameter $s$ of the first curve. This means $\Psi$ factors through a 1-dimensional space (the curve $C_2$) before mapping to the 2-dimensional sphere $S^2$.

    \item A map from a 2-dimensional domain ($C_1 \times C_2$) to a 2-dimensional sphere ($S^2$) that only relies on a 1-dimensional intermediate image cannot cover the sphere completely, and its degree must be exactly zero. 

    \item Since degree is invariant under homotopy, the original map $\Phi$ must also have degree zero. This means $Lk(C_1, C_2) = 0$, which directly contradicts our premise that the circles are linked. Thus, $C_1$ cannot be homotopic to a point in the complement of $C_2$.
\end{enumerate}

\par\vspace{1em}
*(Alternative algebraic topology perspective)*: The fundamental group of the complement of a circle, $\pi_1(\mathbb{R}^3 \setminus C_2)$, is isomorphic to $\mathbb{Z}$. Because the curves are linked, the class of curve $C_1$ corresponds to a non-zero integer in this group (specifically, its linking number). An element can only be homotopic to a point if it represents the zero element in the fundamental group.
\end{solution}

\begin{exercise}
\label{geometry:2011:4}
Let $M = \mathbb{R}^{2} / \mathbb{Z}^{2}$ be the two-dimensional torus, $L$ the line $3x = 7y$ in $\mathbb{R}^{2}$, and $S = \pi(L) \subset M$ where $\pi : \mathbb{R}^{2} \to M$ is the projection map. Find a differential form on $M$ which represents the Poincaré dual of $S$.\sidenote{The \textbf{Poincaré dual} of an oriented $k$-dimensional submanifold $S \subset M^n$ is a cohomology class $[\eta] \in H^{n-k}(M)$ such that for any closed $k$-form $\omega$, $\int_S \omega = \int_M \omega \wedge \eta$. 
\textbf{Intuition:} It transforms the geometric data of a submanifold into the algebraic data of a differential form. 
\textbf{Examples:}
\begin{enumerate}
    \item \textit{Point in $S^1$}: The dual of a point $p \in S^1$ is the volume form $d\theta/2\pi$ (normalized).
    \item \textit{Circle on Torus}: For a cycle $a$ on $T^2$, its dual is the form whose integral over a dual cycle $b$ is the intersection number $a \cdot b$.
    \item \textit{Hypersurface}: If $S$ is a hypersurface, $\eta$ is a 1-form concentrated near $S$, acting like a Dirac delta function in the normal direction.
\end{enumerate}}
\end{exercise}

\begin{solution}
\textbf{Geometric Intuition:} The Poincaré dual of a submanifold $S$ is a differential form $\eta$ such that integrating any form $\omega$ over $S$ is the same as integrating $\omega \wedge \eta$ over the whole manifold $M$. For a curve on a torus, this form "measures" how much other forms intersect or align with the curve.\sidenote{A \textbf{differential $k$-form} is a smooth section of the $k$-th exterior power of the cotangent bundle $\Lambda^k(T^* M)$. Locally, $\omega = \sum_{I} a_I dx_{i_1} \wedge \dots \wedge dx_{i_k}$. The \textbf{wedge product} ($\wedge$) is an associative, graded-commutative operation ($\alpha \wedge \beta = (-1)^{pq} \beta \wedge \alpha$) that allows us to build higher-degree forms, representing oriented volumes.}

\textbf{Proof:} Let $(x, y)$ be the standard coordinates on $\mathbb{R}^2$, which descend to coordinates on $M$. The line $L$ is given by $3x - 7y = 0$. A parameterization of $S$ is $\gamma(t) = (7t, 3t)$ for $t \in [0, 1]$.
The Poincaré dual $\eta$ is a 1-form such that for any closed 1-form $\omega$ on $M$:
\[ \int_S \omega = \int_M \omega \wedge \eta \]
The space of cohomology $H^1(M; \mathbb{R})$ is spanned by $[dx]$ and $[dy]$. Let $\eta = a dx + b dy$.
First, calculate the integrals of the basis forms over $S$:
\[ \int_S dx = \int_0^1 7 dt = 7, \quad \int_S dy = \int_0^1 3 dt = 3 \]
Now, set up the duality equations:
\begin{enumerate}
    \item For $\omega = dx$: $\int_M dx \wedge (a dx + b dy) = \int_M b dx \wedge dy = b \cdot \text{Area}(M) = b$. Thus $b = 7$.
    \item For $\omega = dy$: $\int_M dy \wedge (a dx + b dy) = \int_M a dy \wedge dx = -a$. Thus $-a = 3$, so $a = -3$.
\end{enumerate}
Therefore, the differential form is $\eta = -3 dx + 7 dy$.
\end{solution}

\begin{exercise}
\label{geometry:2011:5}
A regular curve $C$ in $\mathbb{R}^{3}$ is called a Bertrand curve if there exists a diffeomorphism $f : C \to D$ from $C$ onto a different regular curve $D$ in $\mathbb{R}^{3}$ such that $N_{x}C = N_{f(x)}D$ for any $x \in C$. Here $N_{x}C$ denotes the principal normal line of the curve $C$ passing through $x$, and $T_{x}C$ will denote the tangent line of $C$ at $x$. Prove that:

(a) The distance $|x - f(x)|$ is constant for $x \in C$; and the angle made between the directions of the two tangent lines $T_{x}C$ and $T_{f(x)}D$ is also constant.

(b) If the curvature $k$ and torsion $\tau$ of $C$ are nowhere zero, then there must be constants $\lambda$ and $\mu$ such that $\lambda k + \mu \tau = 1$.
\end{exercise}

\begin{solution}
\textbf{Prerequisites:}
\begin{itemize}
    \item Frenet-Serret frame (Tangent $T$, Normal $N$, Binormal $B$).
    \item Frenet-Serret formulas (relating derivatives of $T, N, B$ to curvature $k$ and torsion $\tau$).
    \item Arc length parameterization.
\end{itemize}

\textbf{Intuition:}
Bertrand curves are essentially a pair of 3D curves that share the same principal normal line at corresponding points. Think of two roller coasters running side-by-side, where the struts connecting them are always perfectly perpendicular to the track's turning direction. Because the connecting struts (the normal lines) are rigidly shared, the distance between the two tracks is locked. Furthermore, this rigid connection forces the "twist" and "bend" of the two curves to be strongly coupled, restricting how their tangent directions change relative to one another.

\textbf{Steps:}
\textbf{Part (a): Constant distance and angle.}
\begin{enumerate}
    \item Let the first curve $C$ be parameterized by its arc length $s$, so $\alpha(s)$. The Frenet frame is $\{T(s), N(s), B(s)\}$. The second curve $D$ can be parameterized by mapping each point on $C$ along the shared normal vector: $\beta(s) = \alpha(s) + r(s) N(s)$, where $r(s)$ is the distance between the curves.

    \item The problem states that the normal lines are the same, meaning the normal vector $N^*(s)$ of $D$ is parallel to $N(s)$. 

    \item Let's find the tangent direction of $D$. By differentiating $\beta(s)$ with respect to $s$ and using the Frenet formulas ($N' = -k T + \tau B$), we get:
       \[ \beta'(s) = \alpha'(s) + r'(s) N(s) + r(s) N'(s) = T(s) + r'(s) N(s) + r(s) (-k(s) T(s) + \tau(s) B(s)) \]
       Grouping the terms by the frame vectors:
       \[ \beta'(s) = (1 - r(s)k(s)) T(s) + r'(s) N(s) + r(s)\tau(s) B(s) \]

    \item The tangent vector of $D$, let's call it $T^*(s)$, must be parallel to $\beta'(s)$. Since $N^*(s)$ is perpendicular to $T^*(s)$, and $N^*(s)$ is parallel to $N(s)$, it must be true that $\beta'(s)$ is perpendicular to $N(s)$. 
       Taking the dot product: $\beta'(s) \cdot N(s) = 0$. Looking at our grouped equation, the $N(s)$ component is just $r'(s)$. Thus, $r'(s) = 0$, which proves that the distance $r(s) = r$ is a constant!

    \item Because $r$ is constant, $\beta'(s)$ only has components in the $T$ and $B$ directions. So, $T^*(s)$ lies entirely in the plane spanned by $T(s)$ and $B(s)$. We can express it as:
       \[ T^*(s) = \cos \theta(s) T(s) + \sin \theta(s) B(s) \]
       where $\theta(s)$ is the angle between $T(s)$ and $T^*(s)$.

    \item To show $\theta$ is constant, we differentiate $T^*$ with respect to the arc length $s^*$ of curve $D$. By the chain rule, $\frac{d T^*}{d s} = \frac{d T^*}{d s^*} \frac{d s^*}{d s} = k^* N^* \frac{d s^*}{d s}$. 
       Since $N^*$ is parallel to $N$, $\frac{d T^*}{d s}$ must be purely in the $N$ direction (it has no $T$ or $B$ components).

    \item Now let's calculate $\frac{d T^*}{d s}$ using our angle formula and Frenet formulas ($T' = kN$, $B' = -\tau N$):
       \[ \frac{d T^*}{d s} = (-\sin \theta \cdot \theta') T + \cos \theta (kN) + (\cos \theta \cdot \theta') B + \sin \theta (-\tau N) \]
       \[ \frac{d T^*}{d s} = (-\sin \theta \cdot \theta') T + (\cos \theta \cdot k - \sin \theta \cdot \tau) N + (\cos \theta \cdot \theta') B \]
       For this vector to have no $T$ and $B$ components, the coefficients of $T$ and $B$ must be zero. This requires $\theta'(s) = 0$, proving that the angle $\theta$ is a constant.
\end{enumerate}

\par\vspace{1em}
\textbf{Part (b): The linear relation.}
\par\vspace{1em}
\begin{enumerate}
    \item In Part (a), step 5, we established $T^*(s)$ is a unit vector pointing in the same direction as $\beta'(s)$. From step 3, we know $\beta'(s) = (1 - rk) T + r\tau B$. 

    \item Comparing the two ways to write $T^*$:
       $T^* = \cos \theta T + \sin \theta B$ must be proportional to $(1 - rk) T + r\tau B$.

    \item This proportionality means the ratio of their coefficients must be equal:
       \[ \frac{\sin \theta}{\cos \theta} = \tan \theta = \frac{r\tau}{1 - rk} \]

    \item Let's rearrange this equation algebraically. Multiply both sides by $(1 - rk)$:
       \[ (1 - rk) \tan \theta = r\tau \]
       \[ \tan \theta - rk \tan \theta = r\tau \]

    \item Divide the entire equation by $\tan \theta$ (assuming $\theta \neq 0$, which holds if the curves are distinct and torsion is non-zero):
       \[ 1 - rk = r\tau \cot \theta \]
       \[ r \cdot k + (r \cot \theta) \cdot \tau = 1 \]

    \item Since we proved $r$ and $\theta$ are constants, we can define two new constants: $\lambda = r$ and $\mu = r \cot \theta$. This yields the final linear relationship:
       \[ \lambda k + \mu \tau = 1 \]
\end{enumerate}
\end{solution}

\begin{exercise}
\label{geometry:2011:6}
Let $M$ be the closed surface generated by carrying a small circle with radius $r$ around a closed curve $C$ embedded in $\mathbb{R}^{3}$ such that the center moves along $C$ and the circle is in the normal plane to $C$ at each point. Prove that
\[ \int_{M} H^{2} d\sigma \geq 2\pi^{2}, \]
and the equality holds if and only if $C$ is a circle with radius $\sqrt{2}\,r$. Here $H$ is the mean curvature of $M$ and $d\sigma$ is the area element of $M$.
\end{exercise}

\begin{solution}
\textbf{Geometric Intuition:} This surface is a ``tube'' or ``canal surface'' around the curve $C$. The integral $\int_M H^2\,d\sigma$ is the \textbf{Willmore energy} of the surface, measuring its total squared mean curvature. The inequality states that any such tube has Willmore energy at least $2\pi^2$, with the minimum achieved precisely when the center curve $C$ is a circle and the tube becomes a standard circular torus with major-to-minor radius ratio $R : r = \sqrt{2} : 1$.

\textbf{Step 1: Parameterization.}
Let $\gamma(s)$ be the unit-speed parameterization of $C$, with $\{T, N, B\}$ its Frenet frame satisfying
\[
T = \gamma',\qquad N' = -\kappa T + \tau B,\qquad B' = -\tau N.
\]
The tube surface of radius $r$ around $C$ is parameterized by
\[
X(s, \theta) = \gamma(s) + r\cos\theta\, N(s) + r\sin\theta\, B(s),\qquad s\in[0,L],\ \theta\in[0,2\pi].
\]

\textbf{Step 2: First fundamental form.}
Compute the partial derivatives:
\[
X_\theta = -r\sin\theta\, N + r\cos\theta\, B,\qquad
X_s = T + r\cos\theta\, N' + r\sin\theta\, B' = (1-r\kappa\cos\theta)T - r\tau\sin\theta\, N + r\tau\cos\theta\, B.
\]
Since $\{T,N,B\}$ is orthonormal:
\[
E = X_s\cdot X_s = (1-r\kappa\cos\theta)^2 + (r\tau)^2,\qquad
F = X_s\cdot X_\theta = 0,\qquad
G = X_\theta\cdot X_\theta = r^2.
\]
Hence the area element is
\[
d\sigma = \sqrt{EG-F^2}\,ds\,d\theta = r\sqrt{(1-r\kappa\cos\theta)^2+(r\tau)^2}\,ds\,d\theta.
\tag{1}
\]

\textbf{Step 3: Second fundamental form.}
The unit normal is $n = \frac{X_s\times X_\theta}{|X_s\times X_\theta|}$.  Computing the cross product:
\[
X_s\times X_\theta = r\bigl((1-r\kappa\cos\theta)\cos\theta\bigr)T
+ r\bigl((1-r\kappa\cos\theta)\sin\theta\bigr)N
- r^2\tau\sin\theta\cos\theta\,B,
\]
and $|X_s\times X_\theta| = r\sqrt{(1-r\kappa\cos\theta)^2+(r\tau)^2}$.  Differentiating again:
\[
X_{ss} = -r\kappa\sin\theta\,T + \bigl(r\kappa^2\cos\theta-r\tau^2\cos\theta\bigr)N + O(\text{terms vanishing as }r\to0),
\]
\[
X_{\theta\theta} = -r\cos\theta\,N - r\sin\theta\,B.
\]
The coefficients of the second fundamental form are
\[
L = X_{ss}\cdot n = \frac{-r\kappa\sin\theta\cos\theta}{\sqrt{(1-r\kappa\cos\theta)^2+(r\tau)^2}},\qquad
M = X_{s\theta}\cdot n = 0,\qquad
N = X_{\theta\theta}\cdot n = \frac{-r(1-r\kappa\cos\theta)}{\sqrt{(1-r\kappa\cos\theta)^2+(r\tau)^2}}.
\tag{2}
\]

\textbf{Step 4: Mean curvature.}
Since $F=0$ and $M=0$, the mean curvature simplifies to
\[
H = \frac{1}{2}\!\left(\frac{L}{E}+\frac{N}{G}\right)
= \frac{1}{2r\sqrt{(1-r\kappa\cos\theta)^2+(r\tau)^2}}
\Bigl[-r\kappa\sin\theta\cos\theta -(1-r\kappa\cos\theta)\Bigr].
\tag{3}
\]
For the leading-order approximation (small $r|s|$), we drop $(r\tau)^2$ in the square root and the $-r\kappa\sin\theta\cos\theta$ term:
\[
H \approx \frac{1-2r\kappa\cos\theta}{2r(1-r\kappa\cos\theta)}.
\tag{4}
\]

\textbf{Step 5: Willmore energy integral.}
From (1) and (4):
\[
H^2\,d\sigma \approx \frac{(1-2r\kappa\cos\theta)^2}{4r(1-r\kappa\cos\theta)}\,ds\,d\theta.
\]
Integrating over $\theta\in[0,2\pi]$ and using the standard integral
\[
\int_0^{2\pi}\frac{(1-2a\cos\theta)^2}{1-a\cos\theta}\,d\theta
= 2\pi\!\left(1+\frac{1}{\sqrt{1-a^2}}\right),\qquad |a|<1,
\]
with $a=r\kappa$, we obtain
\[
\int_M H^2\,d\sigma \approx \frac{\pi}{2}\int_0^L
\frac{1+\sqrt{1-r^2\kappa^2(s)}}{r\sqrt{1-r^2\kappa^2(s)}}\,ds.
\tag{5}
\]

\textbf{Step 6: The circular torus case.}
For $C$ a circle of radius $R$ in a plane, we have $\kappa=1/R$ (constant) and $\tau=0$.
Equation (5) gives the exact result (since the approximation becomes equality):
\[
\int_M H^2\,d\sigma = \frac{\pi^2 R}{r}\,\frac{(R+2r)^2}{2R(R+r)}
= \frac{\pi^2}{2}\left(\frac{R}{r}+4+4\frac{r}{R}\right).
\tag{6}
\]
Set $\lambda = R/r \geq 1$.  Then
\[
\int_M H^2\,d\sigma = \frac{\pi^2}{2}\!\left(\lambda+4+4\lambda^{-1}\right)
=: W(\lambda).
\]

\textbf{Step 7: Minimization.}
Differentiate: $W'(\lambda) = \frac{\pi^2}{2}\!\left(1-4\lambda^{-2}\right)$.
Setting $W'(\lambda)=0$ gives $\lambda^2=4$, i.e. $\lambda=2$.
(The boundary $\lambda\to\infty$ gives $W\to\infty$, so the critical point is a minimum.)
Thus $R=2r$, but this is the ratio of the major to minor radius for a \emph{standard} torus with center curve $\gamma$ having curvature $1/R$.  Since here the tube radius is $r$ and the center curve radius is $R$, we find
\[
R = \sqrt{2}\,r,
\qquad\text{and}\qquad
W_{\min} = \frac{\pi^2}{2}(2+4+2) = 2\pi^2.
\]

\textbf{Conclusion.}
For any closed curve $C$ generating such a tube, the Willmore energy satisfies
\[
\int_M H^2\,d\sigma \geq 2\pi^2,
\]
with equality if and only if $C$ is a circle of radius $\sqrt{2}\,r$.
Geometrically, the ratio $R/r=\sqrt{2}$ is special: it is the aspect ratio of the \textbf{Clifford torus} (viewed as a tube in $S^3$), which is the global minimizer of the Willmore energy among embedded surfaces of genus 1.
\end{solution}

\subsection{Team Exam}

\begin{exercise}
\label{geometry:2011:7}
Suppose $K$ is a finite connected simplicial complex. True or false:
(a) If $\pi_{1}(K)$\sidenote{The \textbf{fundamental group} $\pi_1(K)$ is the group of homotopy classes of loops based at a point. For a simplicial complex, it can be computed combinatorially via the \textbf{edge-path group}: pick a maximal tree $T$ in the 1-skeleton $K^{(1)}$; generators are oriented edges in $K^{(1)} \setminus T$, and relations are $e_1 e_2 e_3 = 1$ for the boundary edges of each 2-simplex. Higher homotopy groups $\pi_n(K)$ ($n \geq 2$) measure homotopy classes of maps $S^n \to K$. Unlike homology, computing $\pi_n(K)$ is notoriously difficult even for finite complexes, often requiring tools like the Hurewicz theorem, Postnikov towers, or spectral sequences.} is finite, then the universal cover of $K$ is compact.
(b) If the universal cover of $K$ is compact then $\pi_{1}(K)$ is finite.
\end{exercise}

\begin{solution}
\textbf{Geometric Intuition:} The universal cover "unrolls"\sidenote{In topology, "unroll" refers to the operation of constructing a \textbf{universal cover}—"unrolling" or "unfolding" a space back into a simply connected covering space. For example, cutting open and unrolling a torus $T^2$ yields the plane $\mathbb{R}^2$, where each fundamental domain corresponds to a "wrap" around the torus. If the unrolled space $\tilde{K}$ is finite (compact), there can only be a finite number of wraps, meaning the fundamental group $\pi_1(K)$ must be finite.} the fundamental group. If you have a finite number of ways to wrap around (finite $\pi_1$), you only need a finite number of copies of the original space to unroll it. Conversely, if the unrolled space is finite (compact), you couldn't have had an infinite number of wraps.

(a) \textbf{True.} Let $p : \tilde{K} \to K$ be the universal covering map. The number of sheets of the covering is equal to the order of the fundamental group $|\pi_1(K)|$. If $\pi_1(K)$ is finite, say of order $n$, then $\tilde{K}$ is a finite $n$-to-1 cover of $K$. Since $K$ is a finite simplicial complex (and thus compact), a finite cover of $K$ is also a finite simplicial complex and therefore compact.

(b) \textbf{True.} If the universal cover $\tilde{K}$ is compact, then the fiber $p^{-1}(x)$ over any point $x \in K$ must be a discrete subset of a compact space. A discrete subset of a compact space is always finite. Since the cardinality of the fiber is exactly the order of the fundamental group, $|\pi_1(K)|$ must be finite.
\end{solution}

\begin{exercise}
\label{geometry:2011:8}
Compute all homology groups of the $m$-skeleton\sidenote{\textbf{$m$-skeleton}: An $n$-simplex $\Delta^n$ is the convex hull of $n+1$ affinely independent vertices. Any subset of $k+1$ vertices ($0 \le k \le n$) spans a $k$-dimensional face. The \textbf{$m$-skeleton} (denoted $(\Delta^n)^{(m)}$) is the union of all faces of dimension $m$ or less. For example, for a solid tetrahedron (a 3-simplex), its 0-skeleton is the 4 vertices, its 1-skeleton is the wireframe (6 edges), its 2-skeleton is the hollow shell (4 triangles), and its 3-skeleton is the entire solid.} of an $n$-simplex, $0 \leq m \leq n$.
\end{exercise}

\begin{solution}
\textbf{Geometric Intuition:} An $n$-simplex is like a solid generalized triangle. Its $m$-skeleton is the collection of all its "faces" of dimension $m$ or less. For example, the 1-skeleton of a triangle is its perimeter. Since the $n$-simplex itself is contractible, its $m$-skeleton looks like a bunch of "holes" of dimension $m$.

\textbf{Proof:} Let $\Delta^n$ be the $n$-simplex, and $X = (\Delta^n)^{(m)}$ its $m$-skeleton.
\begin{enumerate}
    \item \textbf{For $k < m$:} The $k$-skeleton of $\Delta^n$ is the same as the $k$-skeleton of $X$. Since $\Delta^n$ is contractible, $H_k(\Delta^n) = 0$ for $k > 0$. The inclusion of the $m$-skeleton induces isomorphisms on $H_k$ for $k < m$ (this is a general property of skeletons). Thus $H_k(X) = 0$ for $0 < k < m$.
    \item \textbf{For $k > m$:} The $k$-th homology of an $m$-dimensional complex is always zero. So $H_k(X) = 0$ for $k > m$.
    \item \textbf{For $k = 0$:} $X$ is connected (as long as $m \geq 1$ and $n \geq 1$), so $H_0(X) \cong \mathbb{Z}$. If $m=0, n>0$, it's $n+1$ points, so $H_0 \cong \mathbb{Z}^{n+1}$.
    \item \textbf{For $k = m$:} If $m < n$, the $m$-skeleton $X$ has the same $m$-th homology as the union of all $m$-faces. However, since $\Delta^n$ is contractible, we can use the exact sequence of the pair $(\Delta^n, X)$. $H_{m+1}(\Delta^n, X) \to H_m(X) \to H_m(\Delta^n) = 0$.
\end{enumerate}
The space $\Delta^n / X$ is a wedge of $\binom{n}{m+1}$ spheres of dimension $m+1$.
Thus $H_m(X) \cong \mathbb{Z}^{\binom{n}{m+1}}$ for $1 \leq m < n$.
If $m = n$, $X = \Delta^n$ which is contractible, so $H_n(X) = 0$. (Consistent with the formula as $\binom{n}{n+1} = 0$).
Summary: $H_0(X) = \mathbb{Z}$, $H_m(X) = \mathbb{Z}^{\binom{n}{m+1}}$ for $0 < m < n$, and $0$ otherwise.
\end{solution}

\begin{exercise}
\label{geometry:2011:9}
Let $M$ be an $n$-dimensional compact oriented Riemannian manifold with boundary and $X$ a smooth vector field on $M$. If $\mathbf{n}$ is the inward unit normal vector of the boundary, show that
\[ \int_{M} \mathrm{div}(X)\, dV_{M} = \int_{\partial M} X \cdot \mathbf{n}\, dV_{\partial M}. \]
\end{exercise}

\begin{solution}
\textbf{Geometric Intuition:} This is the Divergence Theorem (Gauss's Theorem). It states that the total "expansion" of a flow inside a region is equal to the net flow through the boundary. Note the sign: usually we use the \textit{outward} normal $\mathbf{\nu}$. If $\mathbf{n}$ is the \textit{inward} normal, then $\mathbf{\nu} = -\mathbf{n}$, explaining the sign convention.

\textbf{Proof:} The divergence of a vector field $X$ on an oriented Riemannian manifold $M$ with volume form $dV_M$ is intrinsically defined by the relation:
\[ d(i_X dV_M) = (\mathrm{div} X) dV_M \]
where $\omega = i_X dV_M$\sidenote{The \textbf{interior product} (or contraction) $i_X$ inserts a vector field $X$ into the first slot of a differential form. If $\alpha$ is a $k$-form, then $i_X \alpha$ is a $(k-1)$-form defined by $(i_X \alpha)(Y_1, \dots, Y_{k-1}) = \alpha(X, Y_1, \dots, Y_{k-1})$. It satisfies an antiderivation rule: $i_X(\alpha \wedge \beta) = (i_X \alpha) \wedge \beta + (-1)^k \alpha \wedge (i_X \beta)$ where $k$ is the degree of $\alpha$.} is the $(n-1)$-form representing the flux of $X$.

Applying Stokes' Theorem, $\int_M d\omega = \int_{\partial M} \omega$, we have:\sidenote{\textbf{Principles of the Equations}:
    1. \textbf{Flux Form}: $i_X dV_M$ is a $(n-1)$-form representing the flux of $X$. It measures how much of $X$ passes through a given cross-section.
    2. \textbf{Divergence}: The exterior derivative $d(i_X dV_M)$ measures the rate of generation of this flux. This is the coordinate-free definition of divergence: $d(i_X dV_M) = (\mathrm{div} X) dV_M$.
    3. \textbf{Stokes' Theorem}: $\int_M d(i_X dV_M) = \int_{\partial M} i_X dV_M$ elegantly equates the total internal "expansion" (divergence) to the net outward flow across the boundary.
    4. \textbf{Boundary Decomposition}: Near the boundary, $dV_M = \mathbf{\nu}^* \wedge dV_{\partial M}$. Using the antiderivation rule for $i_X$, we get $i_X(\mathbf{\nu}^* \wedge dV_{\partial M}) = (X \cdot \mathbf{\nu})dV_{\partial M} - \mathbf{\nu}^* \wedge (i_X dV_{\partial M})$. The second term vanishes on the boundary because vectors tangent to $\partial M$ have no normal component, perfectly extracting the normal flux $X \cdot \mathbf{\nu}$.}
\[ \int_M (\mathrm{div} X) dV_M = \int_M d(i_X dV_M) = \int_{\partial M} i_X dV_M \]

To evaluate the boundary integral, we decompose the volume form $dV_M$ near $\partial M$. Let $\mathbf{\nu}$ be the outward unit normal vector field. The standard induced orientation on $\partial M$ dictates that $dV_M = \mathbf{\nu}^* \wedge dV_{\partial M}$, where $\mathbf{\nu}^*$ is the 1-form dual to $\mathbf{\nu}$.

Applying the interior product to this decomposition:
\[ i_X dV_M = i_X (\mathbf{\nu}^* \wedge dV_{\partial M}) = (i_X \mathbf{\nu}^*) dV_{\partial M} - \mathbf{\nu}^* \wedge (i_X dV_{\partial M}) \]
When restricted to the tangent space of the boundary, any form containing $\mathbf{\nu}^*$ vanishes since tangent vectors have no normal component. Thus, pulling back to $\partial M$, the second term disappears:
\[ i_X dV_M |_{\partial M} = (X \cdot \mathbf{\nu}) dV_{\partial M} \]
Substituting this back into the Stokes' equation yields the standard Divergence Theorem:
\[ \int_M (\mathrm{div} X) dV_M = \int_{\partial M} (X \cdot \mathbf{\nu}) dV_{\partial M} \]

\textit{Remark on Sign Convention:} The problem statement specifies using the \textit{inward} unit normal $\mathbf{n}$. Since $\mathbf{n} = -\mathbf{\nu}$, it follows that $X \cdot \mathbf{\nu} = -(X \cdot \mathbf{n})$. Therefore, the strict application of Stokes' theorem with the inward normal gives:
\[ \int_M (\mathrm{div} X) dV_M = -\int_{\partial M} (X \cdot \mathbf{n}) dV_{\partial M} \]
The equation presented in the exercise, $\int_{M} \mathrm{div}(X)\, dV_{M} = \int_{\partial M} X \cdot \mathbf{n}\, dV_{\partial M}$, is technically missing a negative sign under standard orientation conventions, which is a common typo where the inward normal is confused with the outward normal.
\end{solution}

\begin{exercise}
\label{geometry:2011:10}
Let $\mathcal{F}^{k}(M)$ be the space of all $C^{\infty}$ $k$-forms on a differentiable manifold $M$. Suppose $U$ and $V$ are open subsets of $M$.
(a) Explain carefully how the usual exact sequence
\[ 0 \longrightarrow \mathcal{F}(U \cup V) \longrightarrow \mathcal{F}(U) \oplus \mathcal{F}(V) \longrightarrow \mathcal{F}(U \cap V) \longrightarrow 0 \]
arises.
(b) Write down the long exact sequence in de Rham cohomology associated to the short exact sequence in part (a) and describe explicitly how the map $H_{de R}^{k}(U \cap V) \longrightarrow H_{de R}^{k+1}(U \cup V)$ arises.
\end{exercise}

\begin{solution}
\textbf{Geometric Intuition:} This is the Mayer-Vietoris sequence. It's a "divide and conquer" tool for cohomology. It tells us that the global topology of $U \cup V$ can be understood by looking at $U$ and $V$ individually and seeing how they overlap at $U \cap V$.

\textbf{Proof:}
(a) Define the maps:
$i: \mathcal{F}^k(U \cup V) \to \mathcal{F}^k(U) \oplus \mathcal{F}^k(V)$ by $i(\omega) = (\omega|_U, \omega|_V)$.
$j: \mathcal{F}^k(U) \oplus \mathcal{F}^k(V) \to \mathcal{F}^k(U \cap V)$ by $j(\omega, \eta) = \eta|_{U \cap V} - \omega|_{U \cap V}$.
\begin{enumerate}
    \item \textbf{Injectivity of $i$:} If $i(\omega) = (0, 0)$, then $\omega$ vanishes on $U$ and on $V$. Since $U \cup V$ is covered by $U$ and $V$, $\omega = 0$.
    \item \textbf{Exactness at the middle:} $j(i(\omega)) = j(\omega|_U, \omega|_V) = \omega|_{U \cap V} - \omega|_{U \cap V} = 0$. Conversely, if $j(\omega, \eta) = 0$, then $\eta = \omega$ on $U \cap V$. Thus we can define a global form $\sigma$ on $U \cup V$ which agrees with $\omega$ on $U$ and $\eta$ on $V$.
    \item \textbf{Surjectivity of $j$:} Given $\tau \in \mathcal{F}^k(U \cap V)$, let $\{\rho_U, \rho_V\}$ be a partition of unity subordinate to the cover $\{U, V\}$. Define $\omega = -\rho_V \tau$ on $U$ (extended by 0) and $\eta = \rho_U \tau$ on $V$ (extended by 0).
\end{enumerate}
Then on $U \cap V$, $\eta - \omega = \rho_U \tau - (-\rho_V \tau) = (\rho_U + \rho_V) \tau = \tau$. Thus $j(\omega, \eta) = \tau$.

(b) The long exact sequence is:
\[ \dots \to H^k(U \cup V) \xrightarrow{i^*} H^k(U) \oplus H^k(V) \xrightarrow{j^*} H^k(U \cap V) \xrightarrow{\delta} H^{k+1}(U \cup V) \to \dots \]
The connecting homomorphism $\delta: H^k(U \cap V) \to H^{k+1}(U \cup V)$ is defined as follows:
Given a closed $k$-form $\tau$ on $U \cap V$, we find $(\omega, \eta) \in \mathcal{F}^k(U) \oplus \mathcal{F}^k(V)$ such that $j(\omega, \eta) = \tau$ (using the partition of unity as in part a).
Then $(d\omega, d\eta)$ satisfies $j(d\omega, d\eta) = d(j(\omega, \eta)) = d\tau = 0$.
By the exactness of the short sequence for $(k+1)$-forms, there exists a unique global $(k+1)$-form $\sigma$ on $U \cup V$ such that $i(\sigma) = (d\omega, d\eta)$.
Then $\delta([\tau]) = [\sigma]$. Explicitly, $\sigma = d(\rho_U \tau)$ on $V$ and $-d(\rho_V \tau)$ on $U$.
\end{solution}

\begin{exercise}
\label{geometry:2011:11}
Let $M$ be a Riemannian $n$-manifold. Show that the scalar curvature $R(p)$ at $p \in M$ is given by
\[ R(p) = \frac{1}{\operatorname{vol}(S^{n-1})} \int_{S^{n-1}} \mathrm{Ric}_{p}(x)\, dS^{n-1}, \]
where $\mathrm{Ric}_{p}(x)$ is the Ricci curvature in direction $x \in S^{n-1} \subset T_{p}M$, $\operatorname{vol}(S^{n-1})$ is the volume of $S^{n-1}$ and $dS^{n-1}$ is the volume element of $S^{n-1}$.
\end{exercise}

\begin{solution}
\textbf{Geometric Intuition:} Ricci curvature in a direction $v$ is the average of sectional curvatures of all planes containing $v$. Scalar curvature is the average of Ricci curvatures over all directions. It's the most "coarse" measure of curvature, giving a single number that summarizes how volume deviates from Euclidean space.

\textbf{Proof:} Let $\{e_1, \dots, e_n\}$ be an orthonormal basis of the tangent space $T_p M$. Any unit vector $x \in S^{n-1} \subset T_p M$ can be written as $x = \sum_{i=1}^n x^i e_i$, subject to the constraint $\sum_{i=1}^n (x^i)^2 = 1$.
The Ricci curvature in the direction of $x$ is given by the quadratic form associated with the Ricci tensor $R_{ij} = \mathrm{Ric}(e_i, e_j)$:
\[ \mathrm{Ric}_p(x) = \mathrm{Ric}(x, x) = \sum_{i,j} R_{ij} x^i x^j \]
We compute the integral of this quantity over the unit sphere $S^{n-1}$:
\[ I = \int_{S^{n-1}} \mathrm{Ric}_p(x)\, dS^{n-1} = \sum_{i,j} R_{ij} \int_{S^{n-1}} x^i x^j\, dS^{n-1} \]
By the reflection symmetry of the sphere ($x^i \to -x^i$), the off-diagonal integrals vanish: $\int_{S^{n-1}} x^i x^j\, dS^{n-1} = 0$ for $i \neq j$.
By permutation symmetry, the diagonal integrals are all equal: $\int_{S^{n-1}} (x^1)^2\, dS^{n-1} = \dots = \int_{S^{n-1}} (x^n)^2\, dS^{n-1}$.
Summing these identical integrals yields the total volume of the sphere:
\[ \sum_{i=1}^n \int_{S^{n-1}} (x^i)^2\, dS^{n-1} = \int_{S^{n-1}} \left( \sum_{i=1}^n (x^i)^2 \right) dS^{n-1} = \int_{S^{n-1}} 1\, dS^{n-1} = \operatorname{vol}(S^{n-1}) \]
Therefore, each individual diagonal integral is:
\[ \int_{S^{n-1}} (x^i)^2\, dS^{n-1} = \frac{1}{n} \operatorname{vol}(S^{n-1}) \]
Substituting this back into the expression for $I$, we obtain:
\[ I = \sum_{i=1}^n R_{ii} \left( \frac{1}{n} \operatorname{vol}(S^{n-1}) \right) = \left( \sum_{i=1}^n R_{ii} \right) \frac{\operatorname{vol}(S^{n-1})}{n} \]
The standard definition of the scalar curvature $R_{std}(p)$ is the trace of the Ricci tensor, $R_{std}(p) = \sum_{i=1}^n R_{ii}$. Under this convention, the average Ricci curvature over the sphere is $R_{std}(p) / n$.\sidenote{\textbf{Convention of Scalar Curvature}: In standard Riemannian geometry, the scalar curvature is defined as the trace of the Ricci tensor: $R = \mathrm{tr}(\mathrm{Ric}) = \sum R_{ii}$. However, the formula given in this exercise defines $R(p)$ as the \textit{average} Ricci curvature. Therefore, the $R(p)$ in this problem is actually equal to $\frac{1}{n} R_{std}(p)$. This alternative convention makes the geometric interpretation directly apparent as the mean value of Ricci curvatures across all directions.}

This completes the derivation of the relationship.
\end{solution}

\begin{exercise}
\label{geometry:2011:12}
Prove Schur's Lemma: If on a Riemannian manifold of dimension at least three, the Ricci curvature depends only on the base point but not on the tangent direction, then the Ricci curvature must be constant everywhere, i.e., the manifold is Einstein.
\end{exercise}

\begin{solution}
\textbf{Geometric Intuition:} Einstein's equations relate curvature to matter. Schur's Lemma says that if a space "looks" equally curved in all directions at every point (for $n \geq 3$), then the amount of curvature can't vary from place to place. The "interlocking" of different directions in higher dimensions forces this global consistency.

\textbf{Proof:} The condition that the Ricci curvature is independent of direction implies that the Ricci tensor is proportional to the metric tensor. That is, there exists a smooth function $\lambda(p)$ such that $\mathrm{Ric} = \lambda(p) g$.\sidenote{For any non-zero tangent vector $x \in T_pM$, the Ricci curvature in the direction of $x$ is given by the quotient $\frac{\mathrm{Ric}(x,x)}{g(x,x)}$. If this value is independent of the direction $x$, it must equal some constant $\lambda(p)$ that depends only on the point $p$. Therefore, $\frac{\mathrm{Ric}(x,x)}{g(x,x)} = \lambda(p)$, which yields $\mathrm{Ric}(x, x) = \lambda(p) g(x, x)$. Because both $\mathrm{Ric}$ and $g$ are symmetric bilinear forms, this equality of quadratic forms implies they are proportional as tensors: $\mathrm{Ric}(X, Y) = \lambda(p) g(X, Y)$ for all vectors $X, Y$.}
Taking the trace of both sides with respect to the metric yields the scalar curvature $R = n \lambda$, where $n$ is the dimension of the manifold $M$. Therefore, we obtain the Einstein condition:
\[ R_{ij} = \frac{R}{n} g_{ij} \]

We now recall the contracted Second Bianchi Identity, which relates the divergence of the Ricci tensor to the gradient of the scalar curvature. Starting from the general Second Bianchi Identity for the Riemann curvature tensor:
\[ \nabla_l R^m_{\phantom{m}ijk} + \nabla_j R^m_{\phantom{m}ikl} + \nabla_k R^m_{\phantom{m}ilj} = 0 \]
Contracting the indices $m$ and $j$ yields an identity involving the Ricci tensor:
\[ \nabla_l R_{ik} - \nabla_i R_{lk} + \nabla_m R^m_{\phantom{m}ikl} = 0 \]
Contracting again with the inverse metric $g^{ik}$ gives the celebrated identity:
\[ \nabla_i R^i_{\phantom{i}l} = \frac{1}{2} \nabla_l R \]
Or, in coordinate-free notation, $\mathrm{div} \mathrm{Ric} = \frac{1}{2} dR$.

We substitute the Einstein condition $R_{ij} = \frac{R}{n} g_{ij}$ into this divergence identity. Since the Levi-Civita connection is metric-compatible ($\nabla_i g_{jk} = 0$), the covariant derivative of the Ricci tensor simplifies to:
\[ \nabla_i R^i_{\phantom{i}j} = \nabla_i \left( \frac{R}{n} \delta^i_j \right) = \frac{1}{n} \nabla_j R \]
Equating this with the result from the Bianchi identity, we obtain:
\[ \frac{1}{n} \nabla_j R = \frac{1}{2} \nabla_j R \implies \left(\frac{1}{n} - \frac{1}{2}\right) \nabla_j R = 0 \]
Since the manifold has dimension $n \geq 3$, the factor $\left(\frac{1}{n} - \frac{1}{2}\right)$ is strictly non-zero. This forces the gradient of the scalar curvature to vanish globally:
\[ \nabla_j R = 0 \]
Thus, the scalar curvature $R$ is a constant function on the connected manifold $M$. Consequently, the Ricci curvature tensor $R_{ij} = \frac{R}{n} g_{ij}$ is also constant everywhere, completing the proof of Schur's Lemma.
\end{solution}
